% ==============================================================================
% Section 2: Related Literature
% ==============================================================================
\section{Related Literature}
\label{sec:literature}

Our paper contributes to several strands of literature. First, we build on the seminal work of \citet{bajari2003deciding} and \citet{porter1993detection} on statistical tests for bid coordination. \citet{bajari2003deciding} propose two testable implications of competitive bidding: exchangeability (all bidders' bids are drawn from the same distribution conditional on observables) and conditional independence (bid residuals within the same auction are uncorrelated). Rejection of either condition is consistent with collusion. The power of these tests, however, depends on the researcher's ability to separate cost-driven bid variation from strategic coordination. We apply the Bajari-Ye framework separately to FL and non-FL firms, exploiting the FL classification as a natural partition of the bidder population. This design sharpens the tests: rather than comparing all bids in ``suspect'' versus ``clean'' tenders (which requires \textit{ex ante} classification of tenders), we compare FL bid residuals to non-FL residuals \textit{within the same tenders}, using common cost shocks as a built-in control for non-strategic correlation.

Second, we contribute to the growing literature on bid-rigging screens. \citet{imhof2019detecting} develops descriptive screens based on bid distribution moments (coefficient of variation, kurtosis, spread) and demonstrates their performance on Swiss procurement data. \citet{imhof2018screening} evaluates screening performance against known cartels. \citet{abrantesmetz2006a} propose a variance screen for detecting collusion. \citet{hüschelrath2014cartel} apply screens to German procurement. More recently, \citet{chassang2022robust} propose robust screens based on bid rankings that do not require parametric assumptions about cost distributions, achieving detection without specifying the competitive model. \citet{conley2016detecting} develop a method for detecting bidder groups in collusive auctions that operates on network structure rather than bid values. We complement these bid-level approaches with a firm-level marker that requires only participation and outcome data---no bid values---and can serve as a computationally inexpensive first-stage filter before applying more data-intensive bid-level screens.

Third, we contribute to the literature on cover bidding specifically. \citet{porter1999ohio} document the role of phantom bidders in Ohio school milk markets, showing that cartel firms submitted complementary bids in non-designated territories. \citet{pesendorfer2000study} identifies collusion through bid rotation patterns in Florida school milk procurement. \citet{asker2010leniency} studies the internal organization of a bidding cartel in stamp auctions, documenting how designated losers submit complementary bids calibrated to lose. \citet{marshall2012economics} provide a comprehensive theoretical treatment of bidding rings, including the coordination technology required to sustain cover bidding across multiple auctions. Our contribution is to propose a systematic, data-driven method for identifying potential cover bidders across all procurement categories simultaneously, and to distinguish empirically between two regimes of cover bidding that have different strategic implications for cartel organization.

Finally, our setting---a large electronic procurement platform in a developing country---contributes to the nascent literature on collusion detection outside advanced economies. Most empirical work on bid rigging uses data from North America, Western Europe, or Japan, where institutional quality is high, enforcement is credible, and bid data are well-curated. Developing countries face a different environment: minimum bidder requirements are strictly enforced (creating incentives for cover bidders), procurement officers may lack the training or resources to detect suspicious patterns, and competition authorities have limited investigative capacity. Brazil's BEC platform, serving S\~{a}o Paulo's state agencies and municipalities, provides a unique laboratory for studying cover bidding at scale in this institutional context. Our sample of 4.5 million tender-items and 40 million bids makes it one of the largest procurement datasets analyzed for collusion screening to date.
